% Report fragment; needs amsmath, amssymb, amsthm, hyperref/xurl.
\section{The literal PA194AX11G construction and its preservation theorem}

The parent is PA194G from the pinned repository commit
\nolinkurl{32d0d8b7d93a63ad4106c18cdb29b37927acdad2}; the composition
uses its explicit binder and the eleven-entry AX protocol
\cite{generalizer,ax11-parent}.
The parent numeral calculus has constructors
\[
p(a,b)=\tfrac12(a+b)(a+b+1)+b,\quad I(a,b)=4p(a,b),\quad
M(a,b)=4p(a,b)+1,
\]
\[
U(a,B)=4p(B,a)+2,\qquad E(a,b)=4p(a,b)+3.
\]
Its axioms, including the Q-oriented universal-distribution head, and the
ordinary PA calibration are retained with attribution in
\path{machines/AX11G/references/parent_generalizer_proof.tex} and the accompanying
calibration texts. The finite Q premise-exchange certificate is also included.
These are dependencies; the present construction does not independently
complete their model-uniform iPRA calibration.

\subsection{One native row, with an observable recurring effect}

Change exactly
\[
 (2066,1):(0,L,302)\quad\longmapsto\quad(1,R,19).
\]
The first action returns from AX. The second enters the existing stack-push
macro and continues into the full MP operand reader. PA194G already has the
no-push successful-MP branch. Thus this one change implements AX11G without
adding controls. There remain 134 fixed controls, 67 raw axiom readers, and
seven simultaneous folds: $134+67-7=194$.

The dense table has working controls $0,\ldots,193$, initial control zero,
and terminal halt label 194, with no outgoing halt rows. Its SHA-256 is
\begin{quote}\small
\nolinkurl{7e1c6442816e5df05e961f2057719252a1976c9a195a786978bf1e2d3c8e466d}.
\end{quote}
The native table starts in 74 and halts in 254. Its exact relabeling is
included. Every working control is reachable in the pointed transition graph.
The tape is two-way infinite, binary, blank zero, with head initially at zero;
all actions write one bit and move one cell.

This change affects recurring commands. Starting with any fixed old store,
its AX stack gains the old gate as one more entry, and a fourth operand is
consumed. That operand can set the arithmetic offset to any natural number.
On the all-zero store and zero arguments, the old gate is zero, but the
parent AX output stack $W$ is positive because it contains $U(0,0)=2$; hence the new stack
$p(0,W)$ is strictly larger than $W$. The zero-pair alias therefore does not
hide the new behavior. The change is also compatible with the G generalizer:
its saved formula is $U(c,h)$, where $h$ is the original cache.

\subsection{Exact source binding}

The source has 3,672 primitive sites. The old AX return at ordinal 3553 is
replaced by a nineteen-site push, with entry 3553 and normal exit 3572.
That exit is the full MP field reader. Every other old source site remains,
with ordinals above 3553 increased by eighteen. All 3,653 exterior resolved
instructions, including their successor indices, agree under this relocation.
No exterior edge enters the interior of the inserted push. The inserted
block is an exact relocation of the old push at ordinal 50, including its
internal successors and sole normal exit.

The source checker first reconstructs the literal PA194G parent and compares
its dense bytes with the pinned parent. It then checks the sole native edit,
the complete source splice, the unchanged arithmetic-box rows, and every
simultaneous fold context. Independently, two different source resolvers
reconstruct the candidate graph from its path-coded source sites. These
checks identify the command actually implemented; a resemblance between
linear lists of primitive instructions would be insufficient.

\subsection{Universal physical contracts and blank initialization}

The construction pass checks all 4,779 branch-sensitive outcomes, over
arbitrary natural register contents and arbitrary permitted tails and gaps.
The independent physical pass derives its four list-sweep induction lemmas
from the literal rows, without treating the arithmetic boxes as oracles.
It checks the same 4,779 outcomes, yielding strictly positive durations
whose minimum constant coefficient is three. The finite checker records
1,083,079 symbolic operations and maximum 379 for one outcome. These are
checker-work counts, not bounds on the unbounded natural register values.

Exactly 58,154 literal transitions from the all-zero tape reach native
control 84 with head at $-2$ and source entry 25. The logical register vector
is $(2,0,0,0,0,0,0,0,0)$, with 26 spare unit blocks. The exact entire finite
support, and hence the blank exterior, is compared with the encoded source
frame. This is a finite startup theorem and supplies no direct evidence of
actual infinite nonhalting.

A finite source execution can therefore be expanded into a finite literal
execution after startup. Positive duration also reflects finite halting:
if a putative literal halt occurred at time $T$ while the source never
halted, the first $T+1$ positive source segments would construct a literal
run beyond that time. Determinism contradicts the absence of an outgoing
transition from the first halt configuration. One need not assume that an
arbitrary intermediate tape can be decoded as a source frame.

\subsection{Precise logical conclusion}

The independently audited command proof is in \path{machines/AX11G/command_proof.tex}.
It supplies total Push and Pop words, all axiom selectors, the dominated
implication selector, a protected-stack MP compiler, and
\[
 \mathsf{Gen}_i=\mathsf{Push};\mathsf S;
     \mathsf A_{i,0,0,0};\mathsf S^{11}.
\]
The offset-setting command required by ordinary AX11 is absent. All these
identities include zero-stack aliases and early halts; the fourth AX field
is accounted for in the scheduler encoding and the all-store soundness proof.

Consequently the new source searches the same numerical-calculus
refutability condition as the G parent. Composing this new preservation
argument with the explicitly retained ordinary PA calibration gives the
same ordinary mathematical calibration to PA refutability for the new table.
The construction and symbolic checks do not themselves establish
\[
 \mathrm{iPRA}\vdash
 \bigl(\neg\operatorname{Halt}_{\mathrm{blank}}(\mathrm{AX11G})
          \leftrightarrow\operatorname{Con}(\mathrm{PA})\bigr).
\]
For that stronger endpoint, one must internalize the proof/code and
physical/scheduler translations with the needed weak-arithmetic induction
bounds, or discharge the repository's all-model $I\Sigma_1$ calibration
predicate for this literal table. A separate report identifies those
remaining premises. Conditional transport from a calibrated numerical
calculus is legitimate; replacing those premises by the fact that the
Python checkers returned PASS is not.

% Independent mathematical audit of the AX11G composition.
\section{Audit of the combined eleven-entry explicit-binder command}
\label{sec:ax11g-audit}

The proposed PA194AX11G construction combines two existing command ideas:
the PA194G cached generalization $U(c,h)$ and the AX11 convention that saves
the guarded gate and falls through to the full MP reader. The interaction
requires an audit because AX now has a fourth input field, and all stack
selectors move by one position. This section verifies the combined command
algebra independently of a long blank-tape simulation.

Write $[a;w]=p(a,w)$, with the first displayed item at the top of a nested
stack. Let $D=d+x$, $H=h+y$, and let
\[
\begin{array}{c|l}
j&q_j\\\hline
0&g(H,D,c),\text{ with failed guard interpreted as zero}\\
1&\Gamma(c)\\
2&L_6^Q(D,c,H)\\
3&L_5(D,H)\\
4&L_4(D,c,H)\\
5&L_3(D,c,H)\\
6&L_2(D,c,H)\\
7&L_1(D,H)\\
8&\rho(D,c,H,d)\\
9&U(c,h).
\end{array}
\]
The last entry uses the \emph{incoming} cache $h$, not the offset-adjusted
number $H=h+y$. All $q_0,\ldots,q_8$ are divisible by four, while
$q_9\equiv2\pmod4$ and $q_9\ge2$.

\subsection{Exact all-store command effect}

The parent AX command leaves $q_0$ on top, pushes
$[q_1,\ldots,q_9,t;w]$, and clears $d,h$. Appending its literal push macro
therefore yields
\[
 (0,[q_0,q_1,\ldots,q_9,t;w],0,0).
\]
The fall-through MP routine reads a genuine new field $z$. It compares
$I(0,z)$ with zero: equality holds exactly at $z=0$, but both branches
have the same visible top zero and the same stack. Hence the exact effect
for \emph{every} old store and every field value is
\[
 \mathsf A_{c,x,y,z}(t,w,d,h)
       =(0,[q_0,q_1,\ldots,q_9,t;w],z,0).
 \tag{G11-A}
\]
Restricting a proof compiler to $z=0$ is valid. Restricting the soundness
argument to $z=0$ would be invalid, since the scheduler enumerates all
four-field command encodings.

The other commands remain
\[
 \mathsf{MP}_x(t,w,d,h)=
 \begin{cases}(v,w,v,0),&I(t,v)=h,\\
              (0,w,v,0),&I(t,v)\ne h,
 \end{cases}\qquad v=d+x,
\]
and
\[
 \mathsf S(t,w,d,h)=
 \begin{cases}\mathsf{halt},&t+h=1,\\
 (h,\operatorname{tl}w,d,\operatorname{hd}w),&t+h\ne1.
 \end{cases}
\]
The sum guard is exact: it does not halt when $t=h=1$.

\subsection{Total stack words}

Put $A_0=\mathsf A_{0,0,0,0}$ and define
\[
 \mathsf{Push}=A_0;\mathsf S^{10};\mathsf{MP}_0,\qquad
 \mathsf{Pop}=A_0;\mathsf S^{12},\qquad
 \mathsf{Discard}=\mathsf{MP}_0.
\]
Iteration of the Shift equation gives
\begin{align*}
 \mathsf A_{c,x,y,0};\mathsf S^{10}(t,w,d,h)
       &=(q_8,[t;w],0,q_9),\\
 \mathsf A_{c,x,y,0};\mathsf S^{11}(t,w,d,h)
       &=(q_9,w,0,t).
\end{align*}
All these Shifts are nonhalting: sums of the relevant residue-zero entries
cannot be one, and a sum containing $q_9\ge2$ cannot be one. Consequently
\begin{align*}
 \mathsf{Push}(t,w,d,h)&=(0,[t;w],0,0),\\
 \mathsf{Pop}(t,[a;w],d,h)&=(t,w,0,a),\\
 \mathsf{Discard}(t,w,0,h)&=(0,w,0,0).
\end{align*}
The last identity requires $d=0$; it is not asserted for arbitrary $d$.
The proof compiler below maintains that premise at every Discard call.
Push's MP comparison fails because implication and universal formulas have
different residues. Pop is total even if $t=1$: its final tested sum
contains $q_9$. The Cantor identities
$p(0,0)=\operatorname{hd}(0)=\operatorname{tl}(0)=0$ are included, so an
empty stack and a collapsed push of zero are covered.

For $0\le j\le8$, from any $(t,w,0,0)$, the word
\[
 \mathsf{Discard};\mathsf A_{A,D,H,0};
       \mathsf S^{j+2};\mathsf{Pop}^{9-j}
\]
returns $(q_j,w,0,0)$. It exposes the selected entry, then consumes the
remaining generated tail and saved zero using total Pop operations. For an
arbitrary dominated leaf with $v\le D$, use instead
\[
 \mathsf{Discard};\mathsf{MP}_v;
 \mathsf A_{A,D-v,H,0};\mathsf S^{10};\mathsf{Pop},
\]
which returns $(\rho(D,A,H,v),w,0,0)$.

\subsection{The direct generalizer and MP composition}

The correct new generalizer is
\[
 \mathsf{Gen}_i=\mathsf{Push};\mathsf S;
                 \mathsf A_{i,0,0,0};\mathsf S^{11}.
\]
From $(a,w,0,0)$, the first two operations give $(0,w,0,a)$, the new AX
stores $U(i,a)$ in position nine, and eleven Shifts return
$(U(i,a),w,0,0)$. No offset-setting command is needed. All the operations
are nonhalting even at $a=1$, because the intermediate tested formulas have
the residue and size properties established above.

For an MP derivation from $a$ and $I(a,b)$, reserve a zero before the two
recursive calls. If neither call halts, the resulting store immediately
before its final suffix is
\[
 (0,[a,I(a,b),0;w],0,0).
\]
Apply $\mathsf S;\mathsf S;\mathsf{MP}_b;\mathsf{Pop}$. The first Shift
is safe. The second halts if $a=1$; otherwise it exposes
$(a,[0;w],0,I(a,b))$. MP succeeds and gives $(b,[0;w],b,0)$; total Pop
returns $(b,w,0,0)$. This covers $a=b=0$ and the zero-stack aliases.

Induction on a finite numerical derivation therefore constructs a word
$C(e)$ that, from every $(t,w,0,0)$, either halts or returns its conclusion
on top while restoring $w$ and clearing the offset and cache. For a refutation,
$A_0;C(e);\mathsf S$ halts from every incoming command store. The leading
$A_0$ clears arbitrary $d,h$ without requiring the stack to be empty.

\subsection{Four-field scheduling and soundness}

With $r$ denoting the encoded continuation, the actual AX encoding is
\[
 \operatorname{enc}(\mathsf A_{c,x,y,z};r)
       =p(c+2,p(y,p(x,p(z,r)))).
\]
The positive continuation of a nonempty word is one plus its command
encoding; a fresh candidate omits that leading one. The unchanged
three-field AX encoding cannot be reused: the MP reader would consume part
of the supposed continuation as its fourth argument. Components of Cantor
pairing do not exceed their containing code, so positive continuations
strictly decrease. Prefixing a universally halting word by
$\mathsf A_{k,0,0,0}$ gives a cofinal family of fresh candidate codes.
Every prefix terminates and leaves a store to which the all-store compiler
applies. Thus carrying the stack between failed trials causes no gap in
the completeness argument.

For soundness, annotate $t,h$ and every live canonical stack entry with a
derivation of its formula in the same numerical calculus. The new AX adds
only its axiom instances, the saved old top and the valid generalization
$U(c,h)$; its new visible top is truth. The offset $z$ needs no theorem
annotation. MP composes derivations on success and writes truth on failure.
Shift transfers the annotations; if $t+h=1$, an annotated code equals one.
No consistency assumption is used in this local preservation argument.

These arguments show that the combined command system searches exactly
the same refutability condition as PA194G, once both are connected to their
literal source programs and physical tables. They justify the new command
composition and a relative mathematical preservation theorem. They do not
complete the separately identified baseline iPRA calibration or constitute
a newly compiled Lean endpoint.
